\documentclass[10pt,a4paper]{amsart}
\usepackage[margin=19mm]{geometry}
\usepackage{amsmath,amssymb,mathtools}
\usepackage[scr=rsfs,cal=euler]{mathalfa}
\usepackage{xfrac}
\usepackage[unicode,hidelinks,bookmarks=false]{hyperref}
\newcommand{\ie}{i.e.\ }
\newcommand{\cf}{cf.\ }
\newcommand{\resp}{respectively\ }
\newcommand{\Subsection}{\subsection}
\providecommand{\nobreakdash}{\nobreak\hbox{-}}
\setlength{\emergencystretch}{2em}
\multlinegap0pt
\usepackage{bm}
\usepackage{bbm}
\usepackage{dsfont}
\usepackage{xspace}
\usepackage{cancel}
\usepackage{mathtools}
\usepackage{amsthm}
\makeatletter
\@ifundefined{theorem}{\newtheorem{theorem}{Theorem}}{}
\@ifundefined{lemma}{\newtheorem{lemma}[theorem]{Lemma}}{}
\makeatother
\newcommand{\NN}{\mathbb{N}}
\newcommand{\RR}{\mathbb{R}}
\newcommand{\SL}{\mathrm{SL}}
\newcommand{\ZZ}{\mathbb{Z}}
\newcommand{\cN}{\mathcal{N}}
\newcommand{\diag}{\mathrm{diag}}
\newcommand{\dist}{\mathrm{dist}}
\title[Simultaneously bounded and dense orbits for commuting Cartan]{Simultaneously bounded and dense orbits for commuting Cartan actions}
\author{Dmitry Kleinbock, Chengyang Wu}
\date{Source: arXiv:2509.05272v1 (CC BY 4.0)}
\begin{document}
\maketitle

\noindent\emph{Source and licence.} Simultaneously bounded and dense orbits for commuting Cartan actions. Version arXiv:2509.05272v1 (CC BY 4.0), distributed under the Creative Commons Attribution 4.0 International licence, as recorded in the arXiv licence metadata for this exact version and shown on the abstract page https://arxiv.org/abs/2509.05272v1. Full rights notice: licenses/arxiv-2509.05272-CC-BY-4.0.txt.\par\smallskip

\noindent\textit{Editorial scope.} Counted unit: the Lemma whose source label is \texttt{L:Isolation} in the article (source0.tex lines 1531--1534, source label \texttt{L:Isolation}), delivered together with the article's own complete original proof of it (source0.tex lines 1535--1644, one \texttt{proof} environment of 110 lines). No mathematics is written, repaired or re-derived: statement and proof are the article's own text line for line. Required context retained uncounted: none. Every definition, display and label of those lines is retained unchanged. Omitted from this package and not redistributed: 1--1530, 1645--2141. Nothing inside a retained excerpt is omitted or altered: every retained body is delivered exactly as the article's lines are written, so no omission receipt of any kind accompanies this package. Citations: the retained excerpts carry no citation command at all, so no bibliography entry of the article is transcribed and no reference is redistributed. Reference adaptation: none was needed and none was made; every reference inside the delivered unit resolves inside it, so no unresolved reference or citation is produced. Printed slips kept literal: no printed word, symbol, hypothesis or number of the counted statement or of its proof was altered. Layout and class: the article's own class is \texttt{amsart} with packages including latexsym, amsfonts, amssymb, mathrsfs, xy, bbm, enumitem, color, mathtools, euscript, tikz, amsmath, amsthm, verbatim; the wrapper is the lane's pinned \texttt{amsart} editorial wrapper at 10pt on A4, which restates the theorem-like environments the retained text needs and adds the article's own macro definitions that the retained text needs (\texttt{\textbackslash NN}, \texttt{\textbackslash RR}, \texttt{\textbackslash SL}, \texttt{\textbackslash ZZ}, \texttt{\textbackslash cN}, \texttt{\textbackslash diag}, \texttt{\textbackslash dist}), with extra wrapper packages (bm, bbm, dsfont, xspace, cancel, mathtools). The delivered numbering is the unit's own. Assets: no retained span contains a figure environment, a graphics-inclusion command, a PDF-page inclusion, a table or a TikZ picture; no image file is copied into this package, the package declares no asset dependency and no image file is redistributed with it. Licence: version arXiv:2509.05272v1 of this article is distributed under the Creative Commons Attribution 4.0 International licence, as recorded in the arXiv licence metadata for this exact version and shown on the abstract page https://arxiv.org/abs/2509.05272v1. A full rights notice is delivered at licenses/arxiv-2509.05272-CC-BY-4.0.txt. Characters: every retained excerpt is pure ASCII, so the delivered engine, XeLaTeX with its default Latin Modern text font, reports no missing character.
\begin{lemma}\label{L:Isolation} Let $d\geq 3$, and $Ay\subseteq X_d$ be a compact $A$-orbit. Then for any compact subset $K\subseteq X_d$, there exists $\epsilon>0$ such that  \begin{equation*}
		0<\dist(x,Ay)<\epsilon\Longrightarrow A^+x\not\subseteq K.
	\end{equation*}
\end{lemma}

\begin{proof}
	Suppose to the contrary, that is, there exists a compact subset $K$ and a sequence $\{x_n\}_{n\in\NN}$ in $X_d$ such that $0<\dist(x_n,Ay)<1/n$ and $A^+x_n\subseteq K$ for any $n\in\NN$. By compactness of $Ay$, we choose and fix an $\epsilon_0>0$ such that the map \begin{equation*}
		U^{+}(\epsilon_0)\times U^-(\epsilon_0)\times Ay\to X_d,\quad (u^+,u^-,ay)\mapsto u^+u^-ay
	\end{equation*}
	is a diffeomorphism onto an open subset $\cN(Ay,\epsilon_0)$ of $X_d$, where 
	\begin{equation*}
		\begin{aligned}
			U^+(\epsilon_0):=\{u\in \SL_d(\RR): u_{ii}=1\,(1\leq i\leq d);\;u_{ij}=0\,(i>j); |u_{ij}|<\epsilon_0\,(i<j)\};\\
			U^-(\epsilon_0):=\{u\in \SL_d(\RR): u_{ii}=1\,(1\leq i\leq d);\;u_{ij}=0\,(i<j); |u_{ij}|<\epsilon_0\,(i>j)\}.
		\end{aligned}
	\end{equation*}
	For $n$ large enough such that $x_n\in \cN(Ay,\epsilon_0/2)$, we write \begin{equation*}
		x_n=u^+_n\cdot u^-_n\cdot y_n, \text{ where }u^+_n\in U^+(\epsilon_0/2), u^-_n\in U^-(\epsilon_0/2), y_n\in Ay.
	\end{equation*}
	By passing to a subsequence we may assume that $\max\{|(u^+_n)_{ij}|:i<j\}$ is always attained by some $(i_0,j_0)$, and without loss of generality that $(u^+_n)_{i_0,j_0}>0$ for all $n\in\NN$. 
	
	Let \begin{equation*}
		\alpha_t:=\diag(e^{(d-i_0)t}I_{i_0},e^{-i_0t}I_{d-i_0}) \in A^{+} \text{ for }t>0,
	\end{equation*}
	and choose $t_n>0$ such that $e^{dt_n}\cdot(u^+_n)_{i_0,j_0}=\epsilon_0/2$. Then \begin{equation*}
		\alpha_{t_n}x_n=(\alpha_{t_n}u_n^+\alpha_{t_n}^{-1})\cdot (\alpha_{t_n}u_n^-\alpha_{t_n}^{-1})\cdot (\alpha_{t_n}y_n),
	\end{equation*}
	where $\alpha_{t_n}u_n^+\alpha_{t_n}^{-1}\in \overline{U^+(\epsilon_0/2)}$ with 
	\begin{equation*}
		|(\alpha_{t_n}u_n^+\alpha_{t_n}^{-1})_{ij}|=\begin{cases}
			e^{dt_n}|(u_n^+)_{ij}|, \text{ if }i\leq i_0<j\\
			|(u_n^+)_{ij}|, \text{ otherwise }\\
		\end{cases}
		\leq\begin{cases}
			\epsilon_0/2, \text{ if }i\leq i_0<j\\
			\epsilon_0/(2e^{dt_n}), \text{ otherwise }\\
		\end{cases},
	\end{equation*}
	$\alpha_{t_n}u_n^-\alpha_{t_n}^{-1}\in U^-(\epsilon_0/2)$ with 
	$|(\alpha_{t_n}u_n^-\alpha_{t_n}^{-1})_{ij}|\leq |(u^-_n)_{ij}|$,
	and  $\alpha_{t_n}y_n\in Ay$. 
	
	
	
	
	Since $u_n^\pm\to 1_G$ as $n\to +\infty$, we see that $\max\{|(u^\pm_n)_{ij}|:i<j\}\to 0$ as $n\to +\infty$. It follows that $t_n\to +\infty$ and that $\alpha_{t_n}u_n^-\alpha_{t_n}^{-1}\to 1_G$. By passing to a further subsequence, we may assume that $\alpha_{t_n}u_n^+\alpha_{t_n}^{-1}\to u^+_{\infty}\in \overline{U^+(\epsilon_0/2)}$ with \begin{equation*}
		|(u_\infty^+)_{ij}|\leq\begin{cases}
			\epsilon_0/2, \text{ if }i\leq i_0<j\\
			0, \text{ otherwise }\\
		\end{cases} \text{when }i<j,\quad 
		(u_\infty^+)_{i_0,j_0}=\epsilon_0/2,
	\end{equation*}
	and $\alpha_{t_n}y_n\to y_{\infty}\in Ay$ as $n\to +\infty$. Therefore, $x_{\infty}:=\lim\limits_{n\to +\infty}\alpha_{t_n}x_n=u^+_{\infty}y_{\infty}$.
	%, and by definition $A^{+}x_\infty\subseteq K$.
	
	For $k\in \{1,2,\cdots,i_0\}$, let \begin{equation*}
		\beta_s^{(k)}:=
		\begin{cases}
			\diag(e^{-(d-k)s}I_{k},e^{ks}I_{d-k}), \text{ if }k<i_0\\ 
			\diag(e^{-(d-1-i_0)s}I_{i_0},e^{(i_0+1)s}I_{j_0-1-i_0},e^{-(d-1-i_0)s}I_1,e^{(i_0+1)s}I_{d-j_0}), \text{ if }k=i_0\\ 
		\end{cases},
	\end{equation*}
	and $\beta_s:=\prod\limits_{k=1}^{i_0}\beta^{(k)}_s$ for $s>0$.
	%On one hand, for each fixed $s>0$, when $n$ is large enough, we have $\beta_s\cdot \alpha_{t_n}\in A^+$, and hence $\beta_s\cdot x_{\infty}=\lim\limits_{n\to +\infty}\left(\beta_s\cdot\alpha_{t_n}\right)x_n\in K$. On the other hand, 
	For each $s>0$, we have $\beta_s\cdot u_\infty^+y_\infty=(\beta_s u_\infty^+\beta_s^{-1})\cdot (\beta_sy_\infty)$, where 
	\begin{equation*}
		|(\beta_su_\infty^+\beta_s^{-1})_{ij}|=\begin{cases}
			\leq e^{-ds}|(u_\infty^+)_{ij}|, \text{ if }i\leq i_0<j \text{ and } (i,j)\neq (i_0,j_0)\\
			|(u_\infty^+)_{ij}|, \text{ if } (i,j)=(i_0,j_0)\\
			0,\text{ otherwise }\\
		\end{cases}
		\text{when }i<j,
	\end{equation*}
	and $\beta_sy_{\infty}\in Ay$. Along a sequence $s_m\to +\infty$, we may assume that 
	%$\beta_{s_m}x_{\infty}\to x_{\infty}'\in K$, 
	$\beta_{s_m}u_\infty^+\beta_{s_m}^{-1}\to I+\frac{\epsilon_0}{2}\cdot E_{i_0,j_0}$, and $\beta_sy_{\infty}\to y_{\infty}'\in Ay$. Write $u_0^+(c):=I+c\cdot E_{i_0,j_0}$. Therefore, it holds that \begin{equation*}
		x_{\infty}':=\lim\limits_{m\to +\infty}\beta_{s_m}\cdot x_{\infty}=u_0^+(\frac{\epsilon_0}{2})\cdot y_{\infty}'.
	\end{equation*}
	
	Next we consider the centralizer of $u_0^+(\frac{\epsilon_0}{2})$ in $A$: it is a $(d-2)$-dimensional subgroup given by \begin{equation*}
		Z_{i_0,j_0}:=\{\diag(a_1,\cdots,a_d)\in A:a_{i_0}=a_{j_0}\}.
	\end{equation*}
	%On one hand, for each fixed $\gamma\in Z_{i_0,j_0}$, On the other hand, 
	By ergodicity it is easy to see that $\overline{Z_{i_0,j_0}\cdot y_{\infty}'}=Ay$. Then we obtain that \begin{equation*}
		\overline{Z_{i_0,j_0}\cdot x_{\infty}'}=u_0^+(\frac{\epsilon_0}{2})\cdot Ay.
	\end{equation*}
	
	Let \begin{equation*}
		\gamma_r:=\diag(e^{(d-i_0)r}I_{i_0},e^{-i_0r}I_{d-i_0})\in A^+ \text{ for } r>0.
	\end{equation*}
	For each $r>0$, we have $\gamma_ru_0^+(\frac{\epsilon_0}{2})\gamma_r^{-1}=u_0^+(e^{dr}\cdot\frac{\epsilon_0}{2})$; in particular, \begin{equation}\label{keyeq}
		\{\gamma_r\cdot \overline{Z_{i_0,j_0}\cdot x_{\infty}'}:r>0\}=\left\{u_0^+(c)\cdot Ay:c>\frac{\epsilon_0}{2}\right\}.
	\end{equation}
	We claim that the left-hand side of (\ref{keyeq}) is bounded but its right-hand side is unbounded, which is a contradiction.
	
	Indeed, for the left-hand side, let us fix any $r>0$, $a\in Z_{i_0,j_0}$, $m\geq 1$. Then when $n$ is large enough, we have $(\gamma_r a \beta_{s_m})\cdot \alpha_{t_n}\in A^+$, and hence 
	\begin{equation*}
		(\gamma_r a \beta_{s_m})\cdot x_{\infty}=\lim\limits_{n\to +\infty}\left((\gamma_r a \beta_{s_m})\cdot\alpha_{t_n}\right)x_n\in K.
	\end{equation*}
	Letting $m\to +\infty$ gives $(\gamma_ra)\cdot x_{\infty}'\in K$. Thus we see that $\{\gamma_r\cdot \overline{Z_{i_0,j_0}\cdot x_{\infty}'}:r>0\}\subseteq K$. For the right-hand side, write $y=h\SL_d(\ZZ)$ and choose a vector $v\in h\ZZ^d$ with $v_{i_0}<0<v_{j_0}$. Let \begin{equation*}
		\delta_t:=\diag(e^{-t}I_{i_0-1},e^{(d-1)t}I_1,e^{-t}I_{d-i_0}) \text{ for }t>0.
	\end{equation*}
	Given any $\epsilon>0$, we may find $t$ large enough such that $|e^{-t}v_i|<\epsilon$ for all $i\neq i_0$, and $c_t:=\frac{e^{(d-1)t}v_{i_0}}{-e^{-t}v_{j_0}}>\frac{\epsilon_0}{2}$. It follows that
	\begin{equation*}
		|(u_0^+(c_t)\cdot \delta_t\cdot v)_i|=\begin{cases}
			e^{-t}|v_i|, \text{ if }i\neq i_0\\
			|e^{(d-1)t}v_{i_0}+c_te^{-t}v_{j_0}|, \text{ if }i=i_0
		\end{cases}
		=\begin{cases}
			<\epsilon, \text{ if }i\neq i_0\\
			0, \text{ if }i=i_0
		\end{cases}.
	\end{equation*}
	This means that $\left\{u_0^+(c)\cdot Ah\ZZ^d:c>\frac{\epsilon_0}{2}\right\}$ contains a nonzero vector with length $<\epsilon$. In view of Mahler's compactness criterion, $\left\{u_0^+(c)\cdot Ay:c>\frac{\epsilon_0}{2}\right\}$ is unbounded. This verifies the claim and completes the proof.
\end{proof}

\end{document}
